% =====================================================================
% VERSION ENSAE, REDUCTION SFE (phase 3, 1er octobre 2026) : annexe A, demonstrations.
% Copie a texte constant de chapitres/15b_demonstrations_cascade.tex.
% Ne changent que les titres et les renvois repares, listes dans
% reduction/phase3.md.
% Lu par main_ensae.tex seul : ne jamais le faire lire par main_institut.tex,
% qui garde chapitres/.
%
% PARTI AU SUPPLEMENT (S5) le 1er octobre 2026 au soir, pour tenir sous 100 pages ; REVENU EN
% ANNEXE A le meme soir, demande de Kelian (« enrichissons le memoire et allons chercher les
% 100 pages ») : c'est l'annexe A du plan cible de la mission. Le supplement passe a S1-S4. Les
% demonstrations sont inchangees ; seule la section sur l'efficience de Shapley est reprise dans
% le registre de la mission (sentence et justification retirees, intertitres en gras supprimes, renvoi
% a « l'annexe des adaptations par pilier » corrige en supplement S1).
% =====================================================================

\chapter{Démonstrations du modèle de cascade}
\label{ann:cascade}

% HARNAIS-HORS-SCRIPT: chapitre de demonstrations de la cascade : memes constantes d'enonce ; les valeurs numeriques correspondantes sont verifiees au chapitre cascade, qui cite ses scripts.

Cette annexe démontre les propriétés énoncées à la section~\ref{sec:formalisation}. Les
notations sont celles du chapitre~\ref{chap:cascade} : $P=\{1,\dots,5\}$, $W\in[0,1]^{P\times P}$
avec $W_{jj}=0$, et la cascade indépendante de la définition~\ref{def:cascade}.

\section*{Forme réduite et descendance espérée (proposition~\ref{prop:reduite})}

\begin{proof}
Par définition du produit matriciel, la somme sur les marches dirigées
$a=v_0\to v_1\to\cdots\to v_n=k$ de $\prod_t W_{v_{t-1}v_t}$ vaut $(W^n)_{ak}$. La série
$\sum_{n\ge0}W^n$ est la série de Neumann de $W$ ; elle converge vers $(I-W)^{-1}$ si et seulement si
$\rho(W)<1$.

Soit $\mathcal C_n(a,k)$ le nombre de chemins vivants de longueur $n$ de $a$ vers $k$, un chemin
étant une marche sans pilier répété. Les arêtes d'un chemin sont distinctes, donc indépendantes, et
un chemin est vivant avec probabilité $\prod_t W_{v_{t-1}v_t}$. Ce n'est pas le cas d'une marche qui
repasse par une arête, celle-ci n'étant tirée qu'une fois. Les chemins formant une partie des
marches, $\mathbb E\,\mathcal C_n(a,k)\le (W^n)_{ak}$. L'indicatrice $\mathbf 1\{k\in S_a\}$ vaut $1$
dès qu'\emph{au moins un} chemin vivant relie $a$ à $k$, donc
$\mathbf 1\{k\in S_a\}\le \sum_{n\ge0}\mathcal C_n(a,k)$ presque sûrement. En prenant l'espérance
puis en sommant sur $k$,
\[
  \mathbb E\,|S_a|=\sum_k \mathbb P(k\in S_a)\le \sum_k\bigl((I-W)^{-1}\bigr)_{ak}
  =\bigl[(I-W)^{-1}\mathbf 1\bigr]_a .
\]
L'inégalité est stricte dès que la série compte des marches qui répètent un pilier, ou deux
chemins distincts vers un même pilier : l'écart mesure ces deux contributions.
\end{proof}

\section*{Sous-criticité par normalisation (proposition~\ref{prop:leontief})}

\begin{proof}
Par construction~\eqref{eq:norm}, $W=(g/s_{\max})\,\mathrm{TRANS}\ge 0$ et, pour toute ligne
$j$, $\sum_k W_{jk}=(g/s_{\max})\sum_k \mathrm{TRANS}_{jk}\le (g/s_{\max})\,s_{\max}=g$, puisque
$s_{\max}$ est la plus grande somme de ligne de $\mathrm{TRANS}$. Pour une matrice à
coefficients positifs, le rayon spectral est majoré par toute norme d'opérateur subordonnée,
en particulier la norme infinie, qui est le maximum des sommes de lignes :
$\rho(W)\le\lVert W\rVert_\infty=\max_j\sum_k W_{jk}\le g\le 1$. Pour $g<1$, on a $\rho(W)<1$ et la
série de Neumann converge donc
et $(I-W)^{-1}=\sum_{n\ge0}W^n$ existe, à coefficients positifs. À la borne $g=1$ du domaine, cette
majoration ne donne que $\rho(W)\le 1$ : la convergence tient alors par la valeur calculée de
$\rho(W)$, strictement inférieure à $1$ (chapitre~\ref{chap:cascade}).
\end{proof}

\section*{Loi exacte de l'ensemble atteint (proposition~\ref{prop:loi})}

\begin{proof}
Fixons $A\ni a$. L'événement $\{S_a=A\}$ est l'intersection de deux événements portant sur des
arêtes disjointes, donc indépendants : (i) aucune arête ne relie $A$ à son complémentaire (sans
quoi $S_a$ dépasserait $A$), et (ii) $a$ atteint tout $A$ par les seules arêtes internes à $A$.
L'événement (i) ne concerne que les arêtes $j\to k$ avec $j\in A$, $k\notin A$, chacune absente
avec probabilité $1-W_{jk}$ indépendamment ; sa probabilité est $B(A)=\prod_{j\in A,k\notin A}(1-W_{jk})$.
L'événement (ii) ne concerne que les arêtes internes à $A$ ; notons $R(A)$ sa probabilité.
Par indépendance, $\mathbb P(S_a=A)=R(A)\,B(A)$.

Pour $R(A)$, on partitionne selon le sous-ensemble $A'$ atteint par les arêtes internes. $A'$
contient $a$, et il est \emph{propre} ($A'\subsetneq A$) si et seulement si $A$ n'est pas atteint
tout entier. Sachant que les arêtes internes atteignent exactement $A'$, aucune arête ne va de $A'$
vers $A\setminus A'$ : cet événement a pour probabilité
$\prod_{j\in A',k\in A\setminus A'}(1-W_{jk})$ et il est indépendant de la réalisation interne à
$A'$. On somme les probabilités de ces sous-ensembles propres, puis on passe au complémentaire :
\[
  R(A)=1-\sum_{\substack{A'\subsetneq A\\ a\in A'}} R(A')\!\!\prod_{j\in A',k\in A\setminus A'}\!\!(1-W_{jk}),
\]
avec $R(\{a\})=1$. La récursion sur les sous-ensembles emboîtés est finie et détermine $R$ de
proche en proche.
\end{proof}

\section*{Monotonie stochastique (proposition~\ref{prop:mono})}

\begin{proof}
On construit un \emph{couplage}. Soit $(U_{jk})$ une famille de variables uniformes sur $[0,1]$
indépendantes, commune aux deux matrices. Pour une matrice $M$, déclarons l'arête $j\to k$
vivante si $U_{jk}<M_{jk}$ ; le graphe obtenu a la bonne loi (arête vivante avec probabilité
$M_{jk}$). Si $W\le W'$ terme à terme, alors $U_{jk}<W_{jk}\Rightarrow U_{jk}<W'_{jk}$ : toute
arête vivante sous $W$ l'est sous $W'$. Le graphe sous $W$ est donc un sous-graphe de celui sous
$W'$, et l'accessibilité étant croissante par ajout d'arêtes, $S_a(W)\subseteq S_a(W')$ presque
sûrement. Pour toute fonctionnelle $\varphi$ croissante pour l'inclusion,
$\varphi(S_a(W))\le\varphi(S_a(W'))$ p.s., d'où la domination stochastique et la monotonie de
l'espérance et des quantiles. Enfin, pour une perturbation antisymétrique $A$, l'entrée $W_{jk}$
croît quand $W_{kj}$ décroît : les deux effets jouent en sens contraire sur $S_a$, la monotonie
en $A_{jk}$ n'a plus lieu, et aucune monotonie coordonnée ne garantit l'optimum aux sommets du pavé $\{|A_{jk}|\le t\,S_{jk}\}$. Leur exploration exhaustive reste un contrôle numérique limité.
\end{proof}

\section*{Dépendance à l'ordre et co-occurrence (propositions~\ref{prop:ordre} et~\ref{prop:nonid})}

\begin{proof}
\emph{Ordre.} Pour une chaîne de deux piliers, $\ell(i,j)=\mathrm{ROOT}_i\,\mathrm{TRANS}_{ij}$ et
$\ell(j,i)=\mathrm{ROOT}_j\,\mathrm{TRANS}_{ji}$ : l'invariance par permutation équivaut à
$\mathrm{ROOT}_i\,\mathrm{TRANS}_{ij}=\mathrm{ROOT}_j\,\mathrm{TRANS}_{ji}$. Une invariance valable
pour toute longueur l'impose en particulier aux chaînes de deux piliers, d'où la nécessité. Au-delà
de deux piliers, la condition ne suffit pas : même avec $\mathrm{TRANS}$ symétrique et des poids
d'amorce égaux, $\ell(1,2,3)$ et $\ell(1,3,2)$ diffèrent dès que
$\mathrm{TRANS}_{12}\mathrm{TRANS}_{23}\neq\mathrm{TRANS}_{13}\mathrm{TRANS}_{32}$. Le barème donne
$\mathrm{ROOT}_1\,\mathrm{TRANS}_{12}\neq\mathrm{ROOT}_2\,\mathrm{TRANS}_{21}$, donc l'ordre de la paire (P1, P2) change la vraisemblance. La différence des criticités arrondies se vérifie séparément dans l'exemple numérique ; elle ne découle pas en général d'une différence de vraisemblance. Un
score additif est une fonction de l'ensemble. Une copule, échangeable ou non, décrit la loi jointe
des indicatrices de défaillance, donc elle aussi une fonction de l'ensemble atteint et non de
l'ordre d'atteinte. Aucun des deux ne reproduit cette dépendance.

\emph{Co-occurrence.} La décomposition $W=\tfrac12(W+W^\top)+\tfrac12(W-W^\top)=S+A$ est
unique car $S$ est symétrique, $A$ antisymétrique, et cette somme directe est celle des
matrices en parties symétrique et antisymétrique.

(i) Pour que $j$ et $k$ soient tous deux atteints depuis une amorce $a\notin\{j,k\}$, il faut au
moins deux arêtes vivantes, événement de probabilité $O(\lVert W\rVert^2)$. Depuis l'amorce $j$,
le pilier $k$ est atteint par l'arête directe avec la probabilité $W_{jk}$, et par tout autre
chemin, de longueur au moins deux, avec une probabilité $O(\lVert W\rVert^2)$ ; symétriquement
depuis l'amorce $k$. En sommant sous l'amorce uniforme, la probabilité de co-occurrence vaut
$r\,(W_{jk}+W_{kj})+O(\lVert W\rVert^2)$, et le terme d'ordre un, $2r\,S_{jk}$, est le même pour
$W$ et pour $W^\top$.

(ii) Deux contre-exemples suffisent. Sous une amorce non uniforme de poids $r_a$, le terme
d'ordre un devient $r_j W_{jk}+r_k W_{kj}$, qui change par transposition dès que $r_j\neq r_k$ et
$W_{jk}\neq W_{kj}$. Sous l'amorce uniforme, prenons trois piliers avec $W_{12}=W_{13}=w>0$ et
toutes les autres entrées nulles. Sous $W$, une amorce en $1$ touche $2$ et $3$ ensemble avec la
probabilité $w^2$, si bien que la paire $\{2,3\}$ co-occurre avec la probabilité $r\,w^2>0$. Sous
$W^\top$, les seules arêtes vont de $2$ et de $3$ vers $1$, et aucun sinistre ne touche $2$ et $3$
ensemble. Les deux matrices ont la même partie symétrique : la co-occurrence d'ordre deux
distingue la fourche du collisionneur, ce qui est la structure en V de l'identification des
graphes causaux \citep{verma1990equivalence}.

(iii) Un changement de signe de $A$ modifie la loi de l'ensemble atteint, donc la loi de la perte
et le SCR en général ; la section~\ref{sec:verif-proprietes} en donne la mesure sur la matrice
publiée.
\end{proof}

\section*{Efficience de la valeur de Shapley}
\label{sec:ann-efficience-shapley}

Le chapitre~\ref{chap:resultats} attribue le surcoût de capital aux piliers par la valeur de
Shapley, dont les parts somment à cent. Cette somme est la propriété d'efficience
\citep{Shapley1953}. Elle permet de présenter l'attribution comme une \emph{partition} du surcoût,
et non comme une pondération arbitraire.

\begin{proposition}[Efficience]
\label{prop:efficience}
Soit $v$ une fonction caractéristique sur $\{1,\dots,n\}$ avec $v(\emptyset)=0$, et
\[
  \varphi_i \;=\; \frac1n \sum_{u \subseteq \{1,\dots,n\}\setminus\{i\}} \binom{n-1}{|u|}^{-1}
                  \bigl[v(u \cup \{i\}) - v(u)\bigr].
\]
Alors $\sum_{i=1}^{n} \varphi_i = v(\{1,\dots,n\})$.
\end{proposition}

\begin{proof}
La démonstration repose entièrement sur la réécriture du poids. Pour $|u|=s$,
\begin{equation}
  \frac1n \binom{n-1}{s}^{-1} \;=\; \frac1n \cdot \frac{s!\,(n-1-s)!}{(n-1)!}
  \;=\; \frac{s!\,(n-s-1)!}{n!},
  \label{eq:poids-shapley}
\end{equation}
de sorte que $\varphi_i$ est exactement la moyenne, sur les $n!$ permutations de
$\{1,\dots,n\}$, de la contribution marginale de $i$ à l'ensemble de ses prédécesseurs. En effet,
un sous-ensemble $u$ de cardinal $s$ ne contenant pas $i$ est l'ensemble des prédécesseurs de $i$
dans exactement $s!\,(n-s-1)!$ permutations : $s!$ ordres pour les éléments de $u$, puis
$(n-s-1)!$ pour ceux qui suivent $i$. En notant $\Pi$ l'ensemble des permutations et $P_i(\pi)$
les prédécesseurs de $i$ dans $\pi$,
\[
  \varphi_i \;=\; \frac{1}{n!} \sum_{\pi \in \Pi}
              \bigl[v\bigl(P_i(\pi) \cup \{i\}\bigr) - v\bigl(P_i(\pi)\bigr)\bigr].
\]
Sommons alors sur $i$ \emph{à permutation fixée}. Si $\pi = (\pi(1),\dots,\pi(n))$, posons $S_0=\emptyset$ et $S_j=\{\pi(1),\dots,\pi(j)\}$ pour $1\le j\le n$. Ainsi $P_{\pi(j)}(\pi)=S_{j-1}$, et la somme des contributions marginales est télescopique :
\[
  \sum_{i=1}^{n} \bigl[v\bigl(P_i(\pi)\cup\{i\}\bigr)-v\bigl(P_i(\pi)\bigr)\bigr]
  \;=\; \sum_{j=1}^{n} \bigl[v(S_j)-v(S_{j-1})\bigr]
  \;=\; v(\{1,\dots,n\}) - v(\emptyset),
\]
et cette valeur ne dépend pas de $\pi$. En moyennant sur les $n!$ permutations,
$\sum_i \varphi_i = v(\{1,\dots,n\})$, puisque $v(\emptyset)=0$.
\end{proof}

L'efficience fait de l'attribution une partition du surcoût, quels que soient l'interaction et son
signe : rien n'est perdu ni compté deux fois. Elle ne fait pas de la part d'un joueur une
prédiction de ce qu'on obtiendrait en l'écartant seul ; cette prédiction est la contribution
marginale, qui ne somme pas au total. Le chapitre~\ref{chap:resultats} tient cette distinction
entre les lectures d'un même canal, et le supplément S1 donne l'identité correspondante pour
chaque colonne. Une part de Shapley est donc une convention d'attribution, pas un budget de
remédiation : deux des quatre canaux du modèle sont bornés et non calibrés.

La réécriture~\eqref{eq:poids-shapley} rend aussi le calcul praticable : elle ramène l'évaluation à
une moyenne sur les permutations, forme sur laquelle reposent les algorithmes d'approximation en
temps polynomial \citep{CastroGomezTejada2009, SongNelsonStaum2016}, nécessaires quand la dimension
croît. Avec cinq piliers, les $5!$ permutations et les $2^5$ valeurs de $v$ sont énumérées en
entier : les valeurs publiées sont exactes au bruit de simulation de $v$ près, sans erreur
d'approximation combinatoire.
